\chaptersection{5}{Risk Assessment}
La prioritizzazione dei rischi viene effettuata mediante un approccio a due livelli, coerente con l'analisi sviluppata nei capitoli precedenti. In una prima fase vengono valutati tutti i failure mode individuati nella FMEA attraverso il Risk Priority Number (RPN), mentre in una fase successiva le minacce di sicurezza selezionate vengono ulteriormente analizzate mediante DREAD.

L'RPN consente di confrontare i failure mode considerando tre fattori: la gravità delle conseguenze (\textit{Severity}), la probabilità relativa di occorrenza (\textit{Occurrence}) e la difficoltà di rilevazione prima che il failure produca effetti (\textit{Detection}). Il valore viene calcolato come:
\[
RPN = S \times O \times D
\]
con ciascun parametro valutato su una scala da 1 a 10. I punteggi non rappresentano probabilità statistiche assolute, ma stime qualitative motivate sulla base dell'architettura, delle cause e degli effetti già individuati nella FMEA.

La valutazione DREAD viene invece applicata successivamente alle minacce di sicurezza ritenute maggiormente rilevanti, considerando \textit{Damage Potential}, \textit{Reproducibility}, \textit{Exploitability}, \textit{Affected Users} e \textit{Discoverability}. In questo modo l'RPN fornisce una prima valutazione della criticità dei failure mode, mentre DREAD permette di approfondire la priorità delle minacce dal punto di vista della sicurezza.


\subsection{Valutazione RPN dei failure mode}
Per ciascun failure mode vengono assegnati i valori di \textit{Severity}, \textit{Occurrence} e \textit{Detection} su una scala da 1 a 10. La valutazione è effettuata in modo comparativo rispetto al contesto specifico di DermaTriage, così da mantenere coerente l'assegnazione dei punteggi tra i diversi failure mode.

In particolare, la \textit{Severity} assume valori più elevati quando il failure può incidere direttamente sulla priorità clinica, compromettere in modo persistente il comportamento della pipeline o causare un'esposizione significativa di dati clinici. L'\textit{Occurrence} cresce per condizioni che possono verificarsi più facilmente durante il normale esercizio del sistema, mentre viene ridotta per
eventi che richiedono condizioni particolari o capacità specifiche da parte di un attaccante. La \textit{Detection} è invece maggiore quando il failure può rimanere nascosto pur lasciando il sistema apparentemente operativo, ed è minore quando l'anomalia risulta immediatamente evidente.

I punteggi ottenuti non rappresentano quindi probabilità assolute, ma una misura relativa della criticità dei failure mode utile per confrontarli e supportare le successive fasi di prioritizzazione.

{
\small
\setstretch{1}
\setlength{\tabcolsep}{4pt}
\renewcommand{\arraystretch}{1.08}

\begin{xltabular}{\textwidth}{
    >{\raggedright\arraybackslash}p{1.2cm}
    X
    >{\centering\arraybackslash}p{0.7cm}
    >{\centering\arraybackslash}p{0.7cm}
    >{\centering\arraybackslash}p{0.7cm}
    >{\centering\arraybackslash}p{1.0cm}
}

\toprule

\textbf{ID} &
\textbf{Failure mode} &
\textbf{S} &
\textbf{O} &
\textbf{D} &
\textbf{RPN} \\

\midrule

\endfirsthead

\toprule

\textbf{ID} &
\textbf{Failure mode} &
\textbf{S} &
\textbf{O} &
\textbf{D} &
\textbf{RPN} \\

\midrule

\endhead

\multicolumn{6}{r}{\footnotesize Continua nella pagina successiva} \\
\endfoot

\endlastfoot


FM01.1 &

Immagine di qualità insufficiente &

7 & 7 & 3 & 147 \\

\midrule


FM01.2 &

Lesione rappresentata in modo incompleto o non adeguato &

8 & 6 & 4 & 192 \\

\midrule


FM01.3 &

Immagine associata al caso clinico errato &

9 & 3 & 7 & 189 \\

\midrule


FM01.4 &

Manipolazione intenzionale dell'immagine &

9 & 4 & 7 & 252 \\

\midrule


FM01.5 &

Immagine non disponibile o non utilizzabile dalla pipeline &

7 & 4 & 2 & 56 \\

\midrule


FM01.6 &

Esposizione non autorizzata dell'immagine dermatologica &

7 & 3 & 7 & 147 \\

\midrule


FM13.1 &

Feedback contenente istruzioni o contenuti avversariali &

9 & 3 & 8 & 216 \\

\midrule


FM13.2 &

Feedback associato al caso clinico errato &

8 & 3 & 7 & 168 \\

\midrule


FM13.3 &

Feedback clinico incompleto o non disponibile &

6 & 4 & 3 & 72 \\

\midrule


FM13.4 &

Feedback clinico alterato o non autentico &

9 & 3 & 8 & 216 \\

\midrule


FM16.1 &

Estrazione inadeguata delle caratteristiche visive &

7 & 5 & 6 & 210 \\

\midrule


FM16.2 &

Classificazione errata dell'urgenza &

9 & 5 & 6 & 270 \\

\midrule


FM16.3 &

Confidence o probabilità non coerenti con la classificazione &

8 & 4 & 8 & 256 \\

\midrule


FM16.4 &

Alterazione o sostituzione non autorizzata dei pesi o del checkpoint &

10 & 2 & 9 & 180 \\

\midrule


FM16.5 &

Output di inferenza incompleto o non disponibile &

7 & 3 & 2 & 42 \\

\midrule


FM20.1 &

Interpretazione errata dell'output di BioMistral &

9 & 3 & 5 & 135 \\

\midrule


FM20.2 &

Determinazione errata della priorità clinica &

10 & 3 & 6 & 180 \\

\midrule


FM20.3 &

Alterazione intenzionale delle regole di mapping &

10 & 3 & 9 & 270 \\

\midrule


FM20.4 &

Produzione incompleta o incoerente del risultato operativo &

8 & 3 & 3 & 72 \\

\midrule


FM21.1 &

Dati del consulto incompleti o incoerenti &

8 & 4 & 4 & 128 \\

\midrule


FM21.2 &

Manipolazione intenzionale dei dati o dei documenti del consulto &

9 & 3 & 8 & 216 \\

\midrule


FM21.3 &

Stato delle informazioni non correttamente aggiornato sulla piattaforma &

7 & 4 & 6 & 168 \\

\midrule


FM21.4 &

Servizi B4 non disponibili o non accessibili &

8 & 4 & 2 & 64 \\

\midrule


FM21.5 &

Esposizione non autorizzata dei dati clinici del consulto &

10 & 3 & 8 & 240 \\

\bottomrule

\caption{Valutazione RPN dei failure mode di DermaTriage}

\label{tab:rpn-failure-modes}\\

\end{xltabular}
}

\subsection{Valutazione DREAD delle minacce}
La valutazione DREAD viene applicata alle minacce individuate mediante il mapping STRIDE-AI, con l'obiettivo di ottenere una prioritizzazione comparativa degli scenari di attacco più rilevanti.

Come primo criterio di selezione viene adottata, ai fini della presente analisi, una soglia pari a
\[
RPN \geq 180,
\]
così da concentrare l'attenzione sui failure mode caratterizzati da una criticità comparativamente elevata.

Il superamento della soglia RPN non è tuttavia sufficiente. La FMEA comprende infatti sia condizioni accidentali sia failure mode riconducibili a un'azione avversariale, mentre DREAD viene utilizzato per la valutazione delle minacce di sicurezza. Vengono quindi mantenuti soltanto i failure mode per i quali il mapping STRIDE-AI consente di individuare una concreta componente intenzionale riconducibile a una categoria STRIDE. Un failure mode può pertanto presentare un RPN elevato senza essere necessariamente sottoposto a DREAD.

Sulla base di tali criteri vengono sottoposti a DREAD i failure mode FM01.3, FM01.4, FM13.1, FM16.4, FM20.3, FM21.2 e FM21.5.

Per ciascuna minaccia vengono valutati \textit{Damage Potential}, \textit{Reproducibility}, \textit{Exploitability}, \textit{Affected Users} e \textit{Discoverability}. Il punteggio ottenuto viene utilizzato come misura comparativa di priorità e non come valore assoluto del rischio, tenendo conto
della componente soggettiva del metodo.

Nell'assegnazione dei punteggi viene mantenuta coerenza con la precedente valutazione FMEA. In particolare, il \textit{Damage Potential} tiene conto della Severity già attribuita al failure mode e dell'eventuale propagazione della compromissione lungo la pipeline. Gli altri fattori vengono stimati considerando le caratteristiche concrete dell'architettura: possibilità di ripetere l'attacco,
prerequisiti tecnici e privilegi necessari, numero di casi o componenti potenzialmente coinvolti e facilità di individuazione della superficie di attacco.

Il punteggio complessivo è calcolato come media dei cinque valori DREAD. Ai soli fini della lettura comparativa dei risultati, nel presente lavoro vengono considerate ad alta priorità le minacce con media almeno pari a 7, a priorità media quelle con valore maggiore o uguale a 5 e inferiore a 7, e a priorità bassa quelle con valore inferiore a 5.

% {
% \small
% \setstretch{1}
% \setlength{\tabcolsep}{3pt}
% \renewcommand{\arraystretch}{1.08}

% \begin{xltabular}{\textwidth}{
%     >{\raggedright\arraybackslash}p{1.15cm}
%     X
%     >{\centering\arraybackslash}p{0.65cm}
%     >{\centering\arraybackslash}p{0.65cm}
%     >{\centering\arraybackslash}p{0.65cm}
%     >{\centering\arraybackslash}p{0.65cm}
%     >{\centering\arraybackslash}p{0.65cm}
%     >{\centering\arraybackslash}p{1.0cm}
%     >{\raggedright\arraybackslash}p{1.25cm}
% }

% \toprule

% \textbf{ID} &
% \textbf{Minaccia valutata} &
% \textbf{D} &
% \textbf{R} &
% \textbf{E} &
% \textbf{A} &
% \textbf{D} &
% \textbf{Media} &
% \textbf{Priorità} \\

% \midrule

% \endfirsthead

% \toprule

% \textbf{ID} &
% \textbf{Minaccia valutata} &
% \textbf{D} &
% \textbf{R} &
% \textbf{E} &
% \textbf{A} &
% \textbf{D} &
% \textbf{Media} &
% \textbf{Priorità} \\

% \midrule

% \endhead

% \multicolumn{9}{r}{\footnotesize Continua nella pagina successiva} \\
% \endfoot

% \endlastfoot


% FM01.3 &
% Falsificazione dell'associazione tra immagine e caso clinico &
% 9 & 5 & 5 & 4 & 5 & 5.6 &
% Media \\

% \midrule

% FM01.4 &
% Manipolazione intenzionale dell'immagine &
% 9 & 7 & 7 & 5 & 6 & 6.8 &
% Media \\

% \midrule

% FM13.1 &
% Inserimento di contenuti avversariali nel feedback clinico &
% 9 & 6 & 5 & 8 & 5 & 6.6 &
% Media \\

% \midrule

% FM16.4 &
% Alterazione o sostituzione non autorizzata del modello &
% 10 & 5 & 4 & 9 & 4 & 6.4 &
% Media \\

% \midrule

% FM20.3 &
% Alterazione intenzionale delle regole di mapping &
% 10 & 8 & 4 & 9 & 5 & 7.2 &
% Alta \\

% \midrule

% FM21.2 &
% Manipolazione dei dati o dei documenti del consulto &
% 9 & 6 & 5 & 4 & 6 & 6.0 &
% Media \\

% \midrule

% FM21.5 &
% Esposizione non autorizzata dei dati clinici &
% 10 & 7 & 6 & 7 & 6 & 7.2 &
% Alta \\

% \bottomrule

% \caption{Valutazione DREAD delle minacce selezionate per DermaTriage.}
% \label{tab:dread-threats}\\

% \end{xltabular}

{
\small
\setstretch{1}
\setlength{\tabcolsep}{3pt}
\renewcommand{\arraystretch}{1.08}

\begin{xltabular}{\textwidth}{
    >{\raggedright\arraybackslash}p{1.25cm}
    X
    >{\centering\arraybackslash}p{0.65cm}
    >{\centering\arraybackslash}p{0.65cm}
    >{\centering\arraybackslash}p{0.65cm}
    >{\centering\arraybackslash}p{0.65cm}
    >{\centering\arraybackslash}p{0.65cm}
    >{\centering\arraybackslash}p{1.2cm}
    >{\centering\arraybackslash}p{1.6cm}
}

\toprule

\textbf{ID} &
\textbf{Minaccia valutata} &
\textbf{D} &
\textbf{R} &
\textbf{E} &
\textbf{A} &
\textbf{D} &
\textbf{Media} &
\textbf{Priorità} \\

\midrule

\endfirsthead

\toprule

\textbf{ID} &
\textbf{Minaccia valutata} &
\textbf{D} &
\textbf{R} &
\textbf{E} &
\textbf{A} &
\textbf{D} &
\textbf{Media} &
\textbf{Priorità} \\

\midrule

\endhead

\multicolumn{9}{r}{\footnotesize Continua nella pagina successiva} \\
\endfoot

\endlastfoot


FM01.3 &
Falsificazione dell'associazione tra immagine e caso clinico &
9 & 5 & 5 & 4 & 5 & 5.6 &
Media \\

\midrule

FM01.4 &
Manipolazione intenzionale dell'immagine &
9 & 8 & 7 & 5 & 6 & 7 &
Alta \\

\midrule

FM13.1 &
Inserimento di contenuti avversariali nel feedback clinico &
9 & 6 & 5 & 8 & 5 & 6.6 &
Media \\

\midrule

FM16.4 &
Alterazione o sostituzione non autorizzata del modello &
10 & 5 & 4 & 9 & 4 & 6.4 &
Media \\

\midrule

FM20.3 &
Alterazione intenzionale delle regole di mapping &
10 & 8 & 4 & 9 & 5 & 7.2 &
Alta \\

\midrule

FM21.2 &
Manipolazione dei dati o dei documenti del consulto &
9 & 6 & 5 & 4 & 6 & 6.0 &
Media \\

\midrule

FM21.5 &
Esposizione non autorizzata dei dati clinici &
10 & 7 & 6 & 7 & 6 & 7.2 &
Alta \\

\bottomrule

\caption{Valutazione DREAD delle minacce selezionate per DermaTriage}
\label{tab:dread-threats}\\

\end{xltabular}
}

I risultati evidenziano tre minacce a priorità alta: FM01.4, FM20.3 e FM21.5. La manipolazione intenzionale dell'immagine risulta particolarmente rilevante perché può essere ripetuta con relativa facilità e influenzare più stadi della pipeline. L'alterazione delle regole dell'Adaptation Layer può invece modificare in modo sistematico la priorità clinica pur lasciando il sistema apparentemente
operativo, mentre l'esposizione non autorizzata dei dati clinici può coinvolgere informazioni sensibili associate a più consulti.

FM13.1 e FM16.4 presentano una criticità prossima alla fascia alta. Nel primo caso, contenuti avversariali inseriti nel feedback possono propagarsi attraverso i successivi meccanismi di aggiornamento; nel secondo, la compromissione del modello può produrre effetti persistenti sull'intera pipeline, pur richiedendo capacità o privilegi maggiori e risultando quindi meno facilmente sfruttabile.

FM01.3 e FM21.2 mantengono infine una priorità media, poiché i relativi effetti sono generalmente più circoscritti al singolo caso clinico rispetto alle minacce in grado di alterare stabilmente il comportamento del sistema o coinvolgere un numero maggiore di consulti.

\subsection{Prioritizzazione finale}
I risultati ottenuti mediante RPN e DREAD vengono sintetizzati attraverso una risk matrix qualitativa basata su due dimensioni: \textit{Impact} e \textit{Likelihood}. L'\textit{Impact} tiene conto principalmente della Severity attribuita nella FMEA e del Damage Potential valutato mediante DREAD,
mentre la \textit{Likelihood} considera congiuntamente Occurrence, Reproducibility ed Exploitability.

La risk matrix viene quindi utilizzata come sintesi delle valutazioni precedenti e non come metodo di scoring indipendente. Per rendere uniforme l'interpretazione dei risultati, la combinazione tra Impact e Likelihood viene ricondotta ai livelli di rischio riportati nella Tabella~\ref{tab:risk-matrix}.

\begin{table}[htbp]
    \centering
    \small
    \setstretch{1}
    \renewcommand{\arraystretch}{1.2}

    \begin{tabular}{
        >{\centering\arraybackslash}p{2.5cm}
        >{\centering\arraybackslash}p{2.3cm}
        >{\centering\arraybackslash}p{2.3cm}
        >{\centering\arraybackslash}p{2.3cm}
    }

        \toprule

        \textbf{Impact / Likelihood} &
        \textbf{Bassa} &
        \textbf{Media} &
        \textbf{Alta} \\

        \midrule

        \textbf{Alto} &
        Medio &
        Alto &
        Molto alto \\

        \midrule

        \textbf{Medio} &
        Basso &
        Medio &
        Alto \\

        \midrule

        \textbf{Basso} &
        Basso &
        Basso &
        Medio \\

        \bottomrule

    \end{tabular}

    \caption{Criterio qualitativo utilizzato nella risk matrix.}
    \label{tab:risk-matrix}
\end{table}

Applicando tale criterio alle minacce già valutate mediante DREAD si ottiene la prioritizzazione finale riportata nella Tabella~\ref{tab:final-risk-prioritization}.

\begin{table}[htbp]
    \centering
    \small
    \setstretch{1}
    \setlength{\tabcolsep}{5pt}
    \renewcommand{\arraystretch}{1.1}

    \begin{tabularx}{\textwidth}{
        >{\raggedright\arraybackslash}p{1.2cm}
        X
        >{\centering\arraybackslash}p{1.7cm}
        >{\centering\arraybackslash}p{1.8cm}
        >{\centering\arraybackslash}p{1.8cm}
    }

        \toprule

        \textbf{ID} &
        \textbf{Minaccia} &
        \textbf{Impact} &
        \textbf{Likelihood} &
        \textbf{Rischio} \\

        \midrule

        FM01.3 &
        Falsificazione dell'associazione tra immagine e caso clinico &
        Alto &
        Media &
        Alto \\

        \midrule

        FM01.4 &
        Manipolazione intenzionale dell'immagine &
        Alto &
        Alta &
        Molto alto \\

        \midrule

        FM13.1 &
        Inserimento di contenuti avversariali nel feedback clinico &
        Alto &
        Media &
        Alto \\

        \midrule

        FM16.4 &
        Alterazione o sostituzione non autorizzata del modello &
        Alto &
        Bassa &
        Medio \\

        \midrule

        FM20.3 &
        Alterazione intenzionale delle regole di mapping &
        Alto &
        Media &
        Alto \\

        \midrule

        FM21.2 &
        Manipolazione dei dati o dei documenti del consulto &
        Alto &
        Media &
        Alto \\

        \midrule

        FM21.5 &
        Esposizione non autorizzata dei dati clinici &
        Alto &
        Media &
        Alto \\

        \bottomrule

    \end{tabularx}

    \caption{Prioritizzazione finale delle minacce mediante risk matrix.}
    \label{tab:final-risk-prioritization}
\end{table}

L'esito della valutazione evidenzia FM01.4 come minaccia a rischio molto alto. FM01.3, FM13.1, FM20.3, FM21.2 e FM21.5 risultano invece a rischio alto, mentre FM16.4 presenta un rischio medio a causa della minore Likelihood, pur mantenendo un impatto elevato.

I risultati ottenuti costituiscono la base per la selezione delle minacce da approfondire nei successivi threat scenario e threat tree.